\scanpage{46}
\mspage{- 18 -}
n\textsuperscript{o} 370

\uline{positive sur} $\handwritten{\mathbb{Q}}_p$ \uline{telle que toute classe suivant} $p^n\handwritten{\mathbb{Z}}_p$ ($n$ \uline{entier} $\handwritten{\geq}0$) \uline{soit de mesure} $p^{-n}$.

La restriction de $\handwritten{\alpha}$ à $\handwritten{\mathbb{Z}}_p$ est évidemment une mesure de Haar sur $\handwritten{\mathbb{Z}}_p$.

6. \uline{Définition locale d'une mesure de Haar}.

PROPOSITION 7.- \uline{Soient} $G$ \uline{un groupe localement compact}, $V$ \uline{une partie ouverte de} $G$,
$\handwritten{\mu}$ \uline{une mesure positive non nulle sur} $V$, \uline{ayant la propriété suivante} : \uline{si} $U$ \uline{est une
partie \unsure{ouverte} de} $V$ \uline{et si} $s\handwritten{\in} G$ \uline{est tel que} $sU\handwritten{\subset} V$, \uline{l'image de la mesure} $\handwritten{\mu}_U$,
\uline{induite par} $\handwritten{\mu}$ \uline{sur} $U$, \uline{par l'homéomorphisme} $\unsure{x|\pi}\to sx$ \uline{de} $U$ \uline{sur} $sU$, \uline{est} $\handwritten{\mu}_{sU}$. \uline{Alors il
existe sur} $G$ \uline{une mesure de Haar à gauche} $\handwritten{\alpha}$ \uline{et une seule induisant} $\handwritten{\mu}$ \uline{sur} $V$.

Pour tout $s\handwritten{\in} G$, soit $\handwritten{\mu}_s$ l'image de $\handwritten{\mu}$ par l'homéomorphisme $x\to sx$ de $V$ sur
$sV$. La restriction de $\handwritten{\mu}_s$ à $V\handwritten{\cap} sV$ est l'image de $\handwritten{\mu}_{s^{-1}(V)\handwritten{\cap} V}$ par la restriction de
$x\to sx$ à $s^{-1}(V)\handwritten{\cap} V$ ; par hypothèse, cette image est $\handwritten{\mu}_{V\handwritten{\cap} sV}$. Par
translation, on en conclut que $\handwritten{\mu}_s$ et $\handwritten{\mu}_t$ ont même restriction à $sV\handwritten{\cap} tV$ quels que
soient $s$, $t$. D'après la prop.1 du chap. III, \S~3, n\textsuperscript{o} 1, il existe donc une mesure $\handwritten{\alpha}$
sur $G$ induisant $\handwritten{\mu}_s$ sur $sV$ quel que soit $s$. Il est clair que $\handwritten{\alpha}$ est l'unique mesure
de Haar à gauche sur $G$ induisant $\handwritten{\mu}$ sur $V$.

COROLLAIRE.- \uline{Soient} $G$, $G'$ \uline{deux groupes localement compacts}, $V$ (resp. $V'$) \uline{un voisinage ouvert de l'élément neutre de} $G$ (resp. $G'$), $\handwritten{\varphi}$ \uline{un isomorphisme local de} $G'$ \uline{à} $G$,
\uline{défini sur} $V'$, \uline{tel que} $\handwritten{\varphi}(V')=V$. \uline{Soient} $\handwritten{\alpha}'$ \uline{une mesure de Haar à gauche sur} $G'$,
\uline{et} $\handwritten{\alpha}'_{V'}$ \uline{sa restriction à} $V'$. \uline{Alors} $\handwritten{\varphi}(\handwritten{\alpha}'_{V'})$ \uline{est la restriction à} $V$ \uline{d'une mesure de
Haar à gauche unique} $\handwritten{\alpha}$ \uline{sur} $G$.

Soit $V_1$ un voisinage ouvert de $e$ dans $G$ tel que $V_1V_1^{-1}\handwritten{\subset} V$. Soit $\handwritten{\mu}$ la restriction de $\handwritten{\varphi}(\handwritten{\alpha}'_{V'})$ à $V_1$. Soient $U$ une partie ouverte de $V_1$ et $s\handwritten{\in} G$ tels que $sU\handwritten{\subset} V_1$.
On a $s\handwritten{\in} V_1V_1^{-1}\handwritten{\subset} V$, donc $s=\handwritten{\varphi}(s')$ avec un $s'\handwritten{\in} V'$. Soit $x\handwritten{\in} U$. On a $x=\handwritten{\varphi}(x')$
avec un $x'\handwritten{\in} V'$, donc $sx=\handwritten{\varphi}(s')\handwritten{\varphi}(x')=\handwritten{\varphi}(s'x')$ puisque $sx\handwritten{\in} sU\handwritten{\subset} V$. Comme les
translations à gauche dans $G'$ conservent $\handwritten{\alpha}'$, on voit que $V_1$ et $\handwritten{\mu}$ satisfont aux
conditions de la prop. 7. Soit $\handwritten{\alpha}$ la mesure de Haar à gauche sur $G$ induisant $\handwritten{\mu}$ sur
% archive sheet number: 45
